\documentclass[11pt]{article}
\usepackage[margin=1in]{geometry}
\usepackage{amsmath,amssymb}
\usepackage{graphicx}
\usepackage{booktabs}
\usepackage[hidelinks]{hyperref}
\usepackage{caption}
\captionsetup{font=small}
\graphicspath{{figures/}}

\title{\vspace{-2.5em}Belief-State Geometry in Transformers:\\
Reproduction and Unsupervised Recovery via Observable Operator Models}
\author{}
\date{}

\begin{document}
\maketitle
\vspace{-3em}

\section{Introduction}

Shai et al.~\cite{shai2024} show that a transformer trained on sequences from a hidden
data-generating process linearly represents, in its residual stream, the geometry of
\emph{belief states}: the Bayesian posteriors over the generator's hidden state given the
observed history. The relevant object is the \emph{mixed-state presentation} (MSP) from
computational mechanics~\cite{crutchfield1989,shalizi2001}: each history $w$ induces a
belief $b(w)\in\Delta^{n-1}$ over the $n$ hidden states, and the set $\{b(w)\}$ forms a
characteristic fractal (or finite point set) in the simplex. An optimally predicting network
must, and does, encode this geometry affinely in its activations.

Section~A reproduces the central supervised results for two processes, Mess3 and RRXOR.
Section~B is a contribution: an \emph{unsupervised} estimator that recovers the belief-state
geometry from activations and the network's own next-token distribution alone --- using no
ground-truth belief labels and no knowledge of the generating process --- built on the theory
of \emph{observable operator models} (OOMs)~\cite{jaeger2000,thon2015}, presented side by side
with the supervised representation. The Appendix studies the generative \emph{operator
rollout}.

\paragraph{Processes and model.} \emph{Mess3} ($x{=}0.05,a{=}0.85$) is a $3$-state, vocab-$3$
process whose MSP is a dense Sierpinski-like fractal. \emph{RRXOR} is a $5$-state, vocab-$2$
synchronising process whose MSP is a finite set of $36$ belief states. We train the paper's
architecture (App.~A.6: $d_{\mathrm{model}}{=}64$, $4$ layers, $1$ head, $d_{\mathrm{mlp}}{=}256$,
ReLU, LayerNorm, SGD) for $10^6$ steps. Both reach their myopic-optimal loss (Mess3
$\approx 0.7935$; RRXOR $\approx 0.5788$). Throughout, regressions and operator fits use
\emph{prefix deduplication} (causal attention makes a prefix's activation identical at every
occurrence) and \emph{$P(w)$-weighting} by prefix probability, following the convention
of~\cite{riechers2025}.

\section{Section A: Supervised reproduction}

For each process we enumerate every positive-probability length-$10$ input, take the
ground-truth belief after each position, regress the (concatenated, per-layer) residual stream
onto the belief, and project the predictions with the same PCA as the ground-truth simplex.

\subsection{Mess3}

The belief geometry appears in the final-layer residual stream as the characteristic fractal
(Figure~\ref{fig:mess3sup}, belief-decode $R^2\approx0.99$). It is not present at
initialisation but \emph{emerges} over training, and is non-trivial: it survives
cross-validation on held-out prefixes (so it is not memorisation) and collapses to the mean
under a belief-label shuffle (so it is not an artefact of projecting onto a fractal-shaped
target). Figure~\ref{fig:mess3fig6} reproduces the paper's Fig.~6 in full.

\begin{figure}[t]
\centering
\includegraphics[width=0.66\textwidth]{fig03_residual_belief_geometry_mess3.png}
\caption{\textbf{Mess3, paper Fig.~5.} Belief geometry recovered from the final-layer residual
stream by linear regression onto belief and projected to the simplex (datashader-style raster).}
\label{fig:mess3sup}
\end{figure}

\begin{figure}[t]
\centering
\includegraphics[width=0.98\textwidth]{fig33_mess3_fig6_supervised.png}
\caption{\textbf{Mess3, paper Fig.~6 (supervised).} \emph{Top:} belief-fractal emergence across
training checkpoints (steps $100$ to $10^6$). \emph{Bottom:} cross-validation (the fractal
persists on held-out prefixes), shuffle control (decode collapses to the mean), and
belief-reconstruction MSE at each checkpoint plus cross-validation and shuffle on the final model.}
\label{fig:mess3fig6}
\end{figure}

\clearpage
\subsection{RRXOR}

RRXOR's belief geometry is the symmetric four-apex diamond of the $36$ MSP states in the
PC1/PC3 plane (left panel of Figure~\ref{fig:rrxorsup}). It is recovered from the concatenated
residual stream (right panel) at $R^2\approx0.87$ (prefix-deduplicated, $P(w)$-weighted),
with the $36$ state centres of mass at $R^2=1.00$. Two further controls match the paper
(Figure~\ref{fig:rrxordiag}): the geometry is spread across \emph{all} layers (every single
layer carries it poorly, the concatenation best), and belief distances are preserved
($R^2\approx0.78$) while next-token distances are not ($R^2\approx0.31$) --- the geometry is not
a next-token readout.

\begin{figure}[t]
\centering
\includegraphics[width=0.82\textwidth]{fig10_rrxor_residual_representation.png}
\caption{\textbf{RRXOR, paper Fig.~7C.} Ground truth (left) and supervised recovery from the
concatenated residual stream (right; prefix-deduplicated, $P(w)$-weighted; small dots =
per-unique-prefix predictions, large = per-state COM, grey rings = ground truth).}
\label{fig:rrxorsup}
\end{figure}

\begin{figure}[t]
\centering
\includegraphics[width=0.98\textwidth]{fig34_rrxor_diagnostics_supervised.png}
\caption{\textbf{RRXOR diagnostics, supervised (paper Fig.~7D--E).} \emph{Left:} per-activation
belief-reconstruction MSE --- RRXOR's geometry is spread across all layers (every single layer
high; the concatenation lowest), while Mess3 carries it in every layer. \emph{Centre/right:}
over the $36$ belief states, the representation preserves belief distances ($R^2\approx0.78$)
but not next-token distances ($R^2\approx0.31$) --- the geometry is more than next-token.}
\label{fig:rrxordiag}
\end{figure}

\clearpage
\section{Section B: Supervised vs.\ unsupervised recovery}

\paragraph{Observable operator models.} A discrete stochastic process over an alphabet $\Sigma$
admits an OOM~\cite{jaeger2000}: a state space $\mathbb{R}^{d}$, one \emph{observable operator}
$\tau_x\in\mathbb{R}^{d\times d}$ per symbol, an initial state $\pi$, and an evaluation
functional $\sigma$, with every word probability
\begin{equation}
\Pr(x_{1}\cdots x_{L}) \;=\; \sigma\,\tau_{x_L}\cdots\tau_{x_1}\,\pi .
\label{eq:oom}
\end{equation}
Hidden Markov models are the special case $\tau_x=T^x$, $\sigma=\mathbf{1}^\top$, with
$T^x_{ij}=\Pr(\text{state }j,\,\text{emit }x\mid\text{state }i)$. The Bayesian \emph{belief}
$b(w)$ is the OOM state after $w$, updated projectively
\begin{equation}
b(wx)=\frac{b(w)\,T^{x}}{b(w)\,T^{x}\mathbf{1}},\qquad
\Pr(x\mid w)=b(w)\,T^{x}\mathbf{1}.
\label{eq:update}
\end{equation}
OOMs, predictive-state representations~\cite{littman2001} and multiplicity automata are unified
by Thon and Jaeger~\cite{thon2015}; the subspace-learning idea below is the one used by Hsu,
Kakade and Zhang for HMMs~\cite{hsu2012}: a process of OOM-dimension $d$ is recovered, up to a
similarity (gauge) transform, from the rank-$d$ factorisation of a past$\times$future statistic.

\paragraph{The observability-OOM estimator.} Section~A establishes that the network encodes
belief affinely, $a(w)=b(w)\Psi$. The projective update~\eqref{eq:update} is non-linear in $a$,
but multiplying through by the normaliser --- which the network \emph{outputs} as its softmax
$\Pr(x\mid w)$ --- linearises it: $a(w)A_x\approx\Pr(x\mid w)\,a(wx)$ with $A_x=\Psi^{+}T^{x}\Psi$.
Let $C:\,a(w)\mapsto\Pr(\cdot\mid w)$ be the one-step observable and $\{G_x\}$ these rescaled
operators (all $P(w)$-weighted ridge regressions). The \emph{observability matrix}
$\mathcal{O}=[\,C\;\;G_xC\;\;G_xG_yC\;\cdots]$ stacks the multi-step observables, each linear in
belief; an SVD of $\mathcal{O}$ returns an orthonormal basis $B$, with the rank giving the model
order $d$. Everything uses only activations, the model softmax and the prefix tree; ground-truth
belief is used solely to \emph{score}. The construction assumes only linearity, so it transfers
unchanged to quantum/post-quantum representations.

The figures below are strict like-for-like: identical pipeline, identical PCA, identical
(label-fit) placement onto the simplex; the \emph{only} difference between the supervised and
unsupervised panels is what feeds the affine readout --- the full residual stream versus the
recovered $d$-dimensional observability subspace.

\subsection{Mess3}

The observability spectrum drops cleanly at $d{=}3$, $\mathrm{eig}(A_x)$ matches
$\mathrm{eig}(T^x)$, and the recovered $3$-D subspace decodes belief at $R^2=0.998$ ---
indistinguishable from the supervised $0.999$ (Figure~\ref{fig:mess3comp}). Mess3 is strongly
observable, so the labels add essentially nothing.

\begin{figure}[t]
\centering
\includegraphics[width=0.96\textwidth]{fig28_mess3_principled.png}
\caption{\textbf{Mess3, supervised vs.\ unsupervised.} Ground truth (left); supervised
full-residual decode (centre, $R^2=0.999$); unsupervised observability-OOM subspace (right,
$R^2=0.998$). Identical pipeline; only the affine-map input differs.}
\label{fig:mess3comp}
\end{figure}

The agreement holds over training and under the same controls as Section~A, with the identical
observability-OOM + $P(w)$ estimator. The belief fractal emerges in lockstep in the supervised
and unsupervised representations across checkpoints (Figure~\ref{fig:mess3progB}), and both
survive cross-validation and collapse under shuffle at comparable reconstruction error
(Figure~\ref{fig:mess3ctrlB}). Both are cross-validated on a stringent $30/70$ split (train
$30\%$, test $70\%$). This is the lowest train fraction the unsupervised method tolerates: its
training unit is a \emph{transition} (a prefix together with all its children), which survives a
random prefix split at only $\sim\!f^4$, so below $\sim\!30\%$ train it runs out of transitions
to estimate the subspace; the supervised decode, whose unit is a single prefix, is unaffected.

\begin{figure}[t]
\centering
\includegraphics[width=0.96\textwidth]{fig06_training_progression_mess3.png}
\caption{\textbf{Mess3 emergence, supervised vs.\ unsupervised.} Belief-fractal emergence across
training checkpoints: top row = supervised residual-stream decode; bottom row = unsupervised
recovery from the raw activations. The geometry sharpens in lockstep in both.}
\label{fig:mess3progB}
\end{figure}

\begin{figure}[t]
\centering
\includegraphics[width=0.98\textwidth]{fig08_unsupervised_controls_mess3.png}
\caption{\textbf{Mess3 controls, supervised vs.\ unsupervised.} Cross-validation, shuffle
control, and per-checkpoint / CV / shuffle MSE bars, for the supervised (top) and unsupervised
(bottom) representations. Both recover the fractal on held-out prefixes and collapse under shuffle.}
\label{fig:mess3ctrlB}
\end{figure}

\subsection{RRXOR}

For RRXOR the unsupervised subspace recovers the diamond but more diffusely
(Figure~\ref{fig:rrxorcomp}): per-(input,position) $R^2=0.65$ versus the supervised $0.91$, and
state-COM geometry $R^2\approx0.82$ versus $1.00$. The gap is genuine, not an estimator
artefact: several independent unsupervised methods converge to a similar state-level ceiling,
the observability spectrum has no clean rank-$5$ elbow, and a control confirms the limit is
structural --- PCA of the $36$ true state-COMs recovers only $R^2=0.35$, so belief is a
\emph{low-variance} feature even among the state centres, and only predictive (not
variance-based) subspaces work. The cause is \emph{weak observability}: RRXOR's belief has
directions that barely change the predicted future (the multi-step version of the $0.33$
next-token result), surfacing only as faint singular values that need a deep horizon ---
which collides with the $n_{\mathrm{ctx}}{=}10$ coverage wall. The supervised regression
sidesteps this with labels; the gap is the ``value of the labels,'' near zero for the strongly
observable Mess3 and large for RRXOR.

\begin{figure}[t]
\centering
\includegraphics[width=0.96\textwidth]{fig26_rrxor_principled.png}
\caption{\textbf{RRXOR, supervised vs.\ unsupervised.} Ground truth (left); supervised
full-residual recovery (centre, per-position $R^2=0.91$); unsupervised $5$-D observability-OOM
subspace (right, per-position $R^2=0.65$, state-COM $R^2\approx0.82$). Identical pipeline; only
the affine-map input differs. Placement onto the simplex uses a label-fit readout in
\emph{both} panels (display only --- the unsupervised content is the subspace).}
\label{fig:rrxorcomp}
\end{figure}

The two RRXOR controls also have unsupervised counterparts (Figure~\ref{fig:rrxordiagB}). The
``belief spread across layers'' result is recovered without labels: each single layer's
observability-OOM subspace decodes belief poorly while the concatenation does best (left panel,
mirroring Figure~\ref{fig:rrxordiag}). The distance control, however, \emph{reverses}: the
recovered representation preserves next-token distances ($R^2\approx0.68$) more than belief
distances ($R^2\approx0.50$), the opposite of the supervised case. This is consistent with the
rest of the section rather than a contradiction --- the observability-OOM subspace is anchored on
the next-token observable $C$, so its dominant geometry is next-token-shaped, and belief is a
low-variance feature within it: linearly recoverable by the readout of
Figure~\ref{fig:rrxorcomp}, but not the leading distance structure. It is a direct manifestation
of RRXOR's weak observability.

\begin{figure}[t]
\centering
\includegraphics[width=0.98\textwidth]{fig35_rrxor_diagnostics_unsupervised.png}
\caption{\textbf{RRXOR diagnostics, unsupervised counterpart.} \emph{Left:} per-activation
belief-decode MSE from each activation's observability-OOM subspace --- ``belief spread across
layers'' is recovered without labels (cf.\ Figure~\ref{fig:rrxordiag}). \emph{Centre/right:}
distance preservation of the recovered representation --- unlike the supervised case it leans
toward next-token ($R^2\approx0.68$) over belief ($R^2\approx0.50$), because the observability
subspace is anchored on the next-token observable.}
\label{fig:rrxordiagB}
\end{figure}

\section{Conclusion}

The supervised results of Shai et al.~\cite{shai2024} reproduce cleanly for both processes, and
the belief-state geometry can moreover be recovered \emph{unsupervised}, from activations and
the model's own softmax, using a classical observable-operator construction: essentially
perfectly for the strongly observable Mess3, and partially --- to a quantifiable ceiling set by
observability --- for RRXOR. The determining factor is process observability, not the
classical-vs-post-quantum distinction, which suggests the natural next experiment is a richly
sampled post-quantum process, where both subspace and operators should come out clean.

\appendix
\section{Operator rollouts}

The subspace recovery above tests ``is belief linearly present?'' A stronger test is the
generative \emph{operator rollout}: discard the per-prefix activations and regenerate every
belief state purely from the learned operators (seeded only at the root,
$s(wx)=s(w)A_x/\,s(w)A_x e$, with $e$ the recovered evaluation functional). We learn the
operators by the same $P(w)$-weighted rescaled least squares, with an ALS operator-consistency
refinement, and apply the identical recipe to the supervised subspace
(Figures~\ref{fig:mess3suproll},~\ref{fig:rrxorsuproll}) and the unsupervised one
(Figures~\ref{fig:mess3unsuproll},~\ref{fig:rrxorunsuproll}).

For \textbf{Mess3} the rollout reproduces the full fractal at $R^2=1.00$ from \emph{both}
subspaces: its operators are diagonalisable strict contractions whose attractor \emph{is} the
belief geometry. For \textbf{RRXOR} the rollout degrades and does not reach the outer apexes
from either subspace (the COM-geom numbers, $\sim0.6$--$0.7$, are flattered by the label-fit
readout; visually the extreme states stay near the centre). This is intrinsic to RRXOR's
operator algebra --- its true belief-update operators are non-diagonalisable ($T^1$ nilpotent,
$T^0$ defective with a triple eigenvalue), so the belief states are transients of a defective
semigroup, not the attractor of a contraction, and iterating estimated operators cannot recover
the weakly-observable extreme states for \emph{any} subspace. Hence for RRXOR the raw subspace
(Section~B), not the rollout, is the trustworthy representation.

\begin{figure}[h]
\centering
\includegraphics[width=0.92\textwidth]{fig29_mess3_supervised_rollout.png}
\caption{\textbf{Mess3, supervised subspace.} Ground truth | supervised decode | operator
rollout from the supervised subspace ($R^2=1.00$).}
\label{fig:mess3suproll}
\end{figure}

\begin{figure}[h]
\centering
\includegraphics[width=0.92\textwidth]{fig23_mess3_unified.png}
\caption{\textbf{Mess3, unsupervised subspace.} Ground truth | unsupervised observability-OOM
subspace | operator rollout ($R^2=1.00$).}
\label{fig:mess3unsuproll}
\end{figure}

\begin{figure}[h]
\centering
\includegraphics[width=0.92\textwidth]{fig29_rrxor_supervised_rollout.png}
\caption{\textbf{RRXOR, supervised subspace.} Ground truth | supervised decode (COM-geom
$0.95$) | operator rollout (COM-geom $0.63$; the extreme states are not reached).}
\label{fig:rrxorsuproll}
\end{figure}

\begin{figure}[h]
\centering
\includegraphics[width=0.92\textwidth]{fig24_rrxor_unified.png}
\caption{\textbf{RRXOR, unsupervised subspace.} Ground truth | unsupervised observability-OOM
subspace (COM-geom $0.66$) | operator rollout (COM-geom $0.74$; metric flattered, the extreme
states stay central).}
\label{fig:rrxorunsuproll}
\end{figure}

\begin{thebibliography}{9}
\small
\bibitem{shai2024} A.~S. Shai, S.~E. Marzen, L.~Teixeira, A.~G. Oldenziel, P.~M. Riechers.
\emph{Transformers represent belief state geometry in their residual stream.} NeurIPS 2024;
arXiv:2405.15943.
\bibitem{riechers2025} P.~M. Riechers, A.~S. Shai, et al. \emph{Neural networks leverage
nominally quantum and post-quantum representations.} 2025. \emph{(verify authors/venue)}
\bibitem{jaeger2000} H.~Jaeger. \emph{Observable operator models for discrete stochastic time
series.} Neural Computation 12(6):1371--1398, 2000.
\bibitem{thon2015} M.~Thon, H.~Jaeger. \emph{Links between multiplicity automata, observable
operator models and predictive state representations --- a unified learning framework.}
JMLR 16:103--147, 2015.
\bibitem{hsu2012} D.~Hsu, S.~M. Kakade, T.~Zhang. \emph{A spectral algorithm for learning
hidden Markov models.} J.\ Computer and System Sciences 78(5):1460--1480, 2012.
\bibitem{littman2001} M.~L. Littman, R.~S. Sutton, S.~Singh. \emph{Predictive representations
of state.} NeurIPS 2001.
\bibitem{crutchfield1989} J.~P. Crutchfield, K.~Young. \emph{Inferring statistical complexity.}
Physical Review Letters 63(2):105--108, 1989.
\bibitem{shalizi2001} C.~R. Shalizi, J.~P. Crutchfield. \emph{Computational mechanics: Pattern
and prediction, structure and simplicity.} J.\ Statistical Physics 104:817--879, 2001.
\end{thebibliography}

\end{document}
